\begin{abstract}
Large language models have demonstrated strong performance on mathematical reasoning and factual knowledge tasks, yet their ability to reason about GPS coordinates and real-world geography remains underexplored. We introduce \textbf{GPSBench}, a comprehensive benchmark of 57,800 samples across 17 tasks for evaluating geographic spatial reasoning in LLMs. Tasks span geometric coordinate operations (distance calculation, bearing computation, coordinate transformation, spherical geometry) and applied reasoning requiring integration of coordinates with world knowledge (coordinate-to-place mapping, spatial relationship inference, geographic pattern recognition). Unlike tool-augmented benchmarks, GPSBench measures intrinsic capabilities---what LLMs know and can compute from parameters alone. Evaluating 14 state-of-the-art LLMs, we find that GPS reasoning is partially solved but fundamentally incomplete, with geometric computations generally more reliable than coordinate-to-place mappings. Geographic knowledge degrades hierarchically---models identify countries well but struggle with city-level localization---yet performance remains robust to coordinate noise, suggesting genuine coordinate understanding rather than memorization. We further demonstrate that GPS augmentation improves downstream geographic benchmarks, and that finetuning induces trade-offs between geometric computation gains and world knowledge degradation.
\end{abstract}
